\subsection{Modularisation}

To get a set of modules for the logical design, first consider the high-level structure of a \ac{cpu} seen in \autoref{fig:cpu_structure} in \autoref{sec:fundamentals_cpu_structure}. The \ac{alu} is an obvious choice for separating into a distinct functional module, as it literally implements an arithmetic function. Likewise, the control unit can be 
split into its own module. The \ac{isa} defines a set of 8 general purpose registers, all of which are homogenous and can therefore be abstracted away into a structure referred to as a register file. This way, the control logic and \ac{io} lines of each individual register can be shared and the entire file can be reasoned about at a higher level. Since there are many types of memory that can be used in an implementation, it makes sense to abstract the specific memory into a generic one, which can then provide a suitable interface to other modules. All memories, including the program counter, program memory, control signal memory and working memory are therefore split into separate modules. Finally, during development, the need for another module for controlling conditional branching became apparent. It also becomes its own module. At this point, an overview of each module, including its function and inputs and outputs is given.

\subsubsection{Arithmetic logic unit}

The \ac{alu} applies an arithmetic function to two binary numbers and outputs another binary number. It also outputs a set of status flags that inform the programmer about certain qualities of the result. This set of flags is described further in \autoref{sec:implementation_isa_memory_model}. Some \ac{alu} functions use the carry status bit from a previous operation as an additional input. Examples of this are the addition and subtraction functions, as well as the bit-shifting operators, in which the carry status bit determines the value of the bits to be inserted. The \ac{alu} must also know which of its functions to apply, therefore the \ac{alu}-opcode must be passed as an input, too. \autoref{tab:implementation_logical_alu_io} shows a list of all inputs and outputs to the \ac{alu} module.

\begin{table}[H]
    \centering
    \begin{tabular}{|l|l|l|} 
        \hline
        \textbf{\textbf{Signal}} & \textbf{Number of Bits} & \textbf{Description}                             \\ 
        \hline
        Opcode                   & 3                       & Selects the arithmetic function to be performed  \\ 
        \hline
        Operand A                & 8                       & The first operand of the arithmetic function     \\ 
        \hline
        Operand B                & 8                       & The second operand of the arithmetic function    \\ 
        \hline
        Carry In                 & 1                       & The carry-out status flag of the last operation  \\
        \hline
    \end{tabular}
    \label{tab:implementation_logical_alu_io}
    \captionof{table}{Input and output signals of the Arithmetic logic unit}
\end{table}

\subsubsection{Register}

Each register stores one word of data. It continuously outputs its current value and may be written to in a synchronous manner. Therefore, it must receive a clock input, as well as another input to be able to select if the register shall be written to upon the next clock cycle. This input is referred to as the ``Write enable'' line. To clear the value and reset the register, a ``Clear'' input is required. In order to read and write data, which may occur simultaneously, two 8 bit wide connections are required. \autoref{tab:implementation_logical_register_io} shows a list of all inputs and outputs to the register module.

\begin{table}
    \centering
    \begin{tabular}{|l|l|l|} 
        \hline
        \textbf{\textbf{Signal}} & \textbf{Number of Bits} & \textbf{Description}                               \\ 
        \hline
        Data In                  & 8                       & Data to be written                                 \\ 
        \hline
        Data Out                 & 8                       & Readout of the currently stored data               \\ 
        \hline
        Clock                    & 1                       & Controls timing of data writing                    \\ 
        \hline
        Write Enable             & 1                       & Control if data is written or not                  \\ 
        \hline
        Clear                    & 1                       & Asynchronously writes the value 0 to the register  \\
        \hline
    \end{tabular}
    \label{tab:implementation_logical_register_io}
    \captionof{table}{Input and output signals of an individual register}
\end{table}

\subsubsection{Register file}

Since it is not feasible for an implementation to provide a set of control lines and \ac{io} lines for each individual register, a sharing mechanism is required. As there are 8 registers in the register file, $log_2(8) = 3$ bits are required to select a single one. From the instruction encoding described in \autoref{sec:implementation_instruction_encoding}, it is clear that an instruction can simultaneously read from two registers and write to another one. Therefore, each of these accesses requires its own selection input, totaling 9 bits. As only one register is written to at a time, all registers can share their write-enable, clock and data inputs. The register file module has only one of each and distributes it to all its register submodules. Similarly, all registers can share their ``clear'' signal, since no instruction defined in the \ac{isa} requires clearing of a single register. If needed, a programmer can load the constant value 0 from register \texttt{R0} into the desired register to clear it. In addition to general purpose registers, the flags register can be written to simultaneously. Therefore, two sets of write enable and data inputs are required. \autoref{tab:implementation_logical_register_file_io} shows a list of inputs and outputs of the register file module.

\begin{table}
    \centering
    \begin{tabular}{|l|l|l|} 
        \hline
        \textbf{\textbf{Signal}} & \textbf{Number of Bits} & \textbf{Description}                                \\ 
        \hline
        Read Select 0            & 3                       & Selects the first register to read from             \\ 
        \hline
        Read Select 1            & 3                       & Selects the second register to read from            \\ 
        \hline
        Write Select             & 3                       & Selects the register to write to                    \\ 
        \hline
        Clear                    & 1                       & Asynchronously writes the value 0 to all registers  \\ 
        \hline
        Clock                    & 1                       & Controls timing of data writing                     \\ 
        \hline
        Write General Purpose    & 1                       & Control if general purpose data is written or not   \\ 
        \hline
        Write Flag               & 1                       & Control if flag data is written or not              \\ 
        \hline
        Data In General Purpose  & 8                       & Data to be written to write-selected register       \\ 
        \hline
        Data In Flag             & 8                       & Data to be written to flags register                \\
        \hline
    \end{tabular}
    \label{tab:implementation_logical_register_file_io}
    \captionof{table}{Input and output signals of the register file}
\end{table}

\subsubsection{Program counter}
\subsubsection{Memory}